\section{Introduction}

End-to-end testing is a cornerstone of modern continuous delivery pipelines. By simulating real user interactions in a browser, E2E tests validate application behaviour across the full stack without mocking, providing a high-confidence quality gate before every deployment \cite{humble2010continuousdelivery}. Google Chrome in headless mode is the established standard for this purpose: through the Chrome DevTools Protocol (CDP), it is controlled by frameworks such as Playwright \cite{playwright2024} without requiring a display. Its principal drawbacks in CI environments are a large binary footprint (exceeding 100 MB) and significant per-instance memory consumption, which translate into longer pipeline durations and higher compute costs on resource-constrained cloud runners \cite{laukkanen2017problems}.

Lightpanda\footnote{\url{https://lightpanda.io}} is a recently introduced headless browser implemented in Zig, targeting cloud and CI workloads. Its homepage advertises execution time \textbf{9$\times$ faster} and memory usage \textbf{16$\times$ lower} than Chrome, citing a benchmark across 933 ``real web pages.'' It also exposes a CDP endpoint, which in principle allows existing Playwright test suites to connect without modification. These claims motivate a natural question: \emph{Does Lightpanda actually work as a drop-in replacement for Chrome Headless in production E2E test pipelines?}

This paper answers that question through a three-stage empirical study of increasing realism. Stage 1 evaluates raw performance on a controlled local benchmark. Stage 2 exercises Lightpanda in a real GitHub Actions CI pipeline with a simple test application. Stage 3—the most critical stage—runs a comprehensive Playwright suite of 40 tests against a realistic multi-page web application, revealing where Lightpanda's CDP and JavaScript engine coverage breaks down in practice. We additionally examine the validity of Lightpanda's own published benchmark methodology.

The key finding is that while Lightpanda's performance advantages are real and reproducible in synthetic settings, they are rendered moot by severe compatibility failures in real application testing: 30 of 40 tests fail on Lightpanda while passing on Chrome. Lightpanda is not yet production-ready as a Chrome Headless replacement.
